%----------
%	RECUPERACIÓN SOBRE GUÍAS EXPERTAS (RAG)
%----------

El codebook del proyecto se inspira en guías institucionales de lenguaje no
sexista (Sección~\ref{sec:experts}). Esas guías cumplen dos papeles distintos en
el sistema, y conviene no confundirlos.

Como \textbf{fuente}, son el material del que se destila la metodología de cada
variable. Ese uso es indirecto: la analista y quien redacta las skills se apoyan
en ellas, pero el modelo no las lee.

Como \textbf{herramienta en vivo}, en cambio, el agente puede consultarlas
durante la inferencia. La acción \texttt{CONSULTAR\_GUIA} implementa una
recuperación (\textit{Retrieval-Augmented Generation}, RAG por sus siglas en inglés)~\cite{lewis2020rag} sobre el texto real
de las guías. La implementación es deliberadamente ligera: un índice TF-IDF en
memoria sobre los pasajes de las guías, sin dependencias externas. Ante una
consulta, devuelve los pasajes más relevantes con su cita de fuente y sección.

Esta distinción importa para la honestidad del trabajo. Sin la herramienta, la
conexión con las guías sería solo conceptual, y afirmar que \textit{``el agente se apoya en las guías expertas''} sería inexacto. Con \texttt{CONSULTAR\_GUIA}, esa afirmación
se sostiene: el agente recupera texto literal de las fuentes en tiempo de
inferencia y puede citarlo en su justificación.

La recuperación se basa en coincidencia léxica (TF-IDF)~\cite{sparckjones1972}, no semántica. Funciona
bien para términos precisos, como \textit{masculino genérico} o \textit{feminización de
cargos}, pero es sensible a sinónimos y paráfrasis. Una recuperación por
\textit{embeddings} mejoraría esa robustez y queda como línea de mejora. Para el
alcance de este trabajo, el índice léxico es suficiente. Cuestión distinta es si esa recuperación mejora el acuerdo final: su
efecto se evalúa por separado en la Sección~\ref{subsec:iris_descomposicion}, y
resulta, en la práctica, prescindible.
